\section{Calibration details}
\label{app:calib}

This appendix collects the development ensembles of Section~\ref{sec:calib}. They used independent seeds and
smaller samples than the production ensembles of Section~\ref{sec:results}, which supersede them wherever the two
overlap.

\subsection{Procedure}
\label{app:calib-procedure}

\paragraph{Transplants.}
The transplant patterns O12 and O12b are ensemble experiments.
\begin{itemize}
  \item From each replicate colony, four reproductive males and four reproductive females are moved to an
    empty lattice: either from the middle of the growth window (control, phase~B) or at
    $t=1.75\,t_{\rm pk}$ (mid-phase~D).
  \item In O12b, four phase-D animals of one sex are paired with four healthy animals ($s=1$) of the other.
  \item A transplant ``founds'' a colony if it reaches 80 individuals within 150 weeks.
\end{itemize}

\paragraph{Ensembles.}
The calibration proceeded in stages:
\begin{enumerate}
  \item screening of rule families: 23 points $\times$ 10 replicates, then 24 points (initial condition
    $\times\,\alpha\times\kappa$) $\times$ 10 replicates;
  \item ablations of the selected rules: 10 points $\times$ 40 replicates;
  \item one-at-a-time sensitivity around the selected point: 17 points $\times$ 30 replicates;
  \item final ensembles: \textsf{C0} with 120 replicates, 100 transplants and an $\alpha=0$ control;
    \textsf{C1} and \textsf{C2} with 60 replicates and 60 transplants each;
  \item exploratory tests of the founder protocol: 30--60 replicates per point.
\end{enumerate}
Further details:
\begin{itemize}
  \item replicates use independent seeds, \texttt{SeedSequence([base\_seed, point, replicate])};
  \item the horizon is $T=600$ weeks, and every run of the final ensembles of \textsf{C0} and \textsf{C1}
    went extinct before week 220;
  \item a run of \textsf{C0} costs about 1~s of CPU;
  \item the whole calibration used about 45~min on four cores.
\end{itemize}

\subsection{Development ensembles of the candidates}
\label{app:calib-dev}

Table~\ref{tab:calib-dev} reports every pattern in the development ensembles of \textsf{C0}, \textsf{C1} and
\textsf{C2}. Two development estimates were optimistic: O9 in \textsf{C1} (0.57 here, 0.45 in production) and
O17 in \textsf{C2} (0.95 here, 0.83 in production, where 7.5\% of the colonies recover natality).

\begin{table}[t]
  \centering
  \caption{Development ensembles of \textsf{C0}, \textsf{C1} and \textsf{C2}: pattern pass fractions (fraction of
  PASS runs, 95\% Wilson interval; superseded by Table~\ref{tab:res-production}). $N_0=400$, $T=600$ weeks. Transplants: $P_T(\cdot)$ is the fraction of transplants that found a
  colony. ``Persistent O13'' requires the isolated class to stay isolated for $\ge4$ consecutive weeks
  ($\Phi_{\rm BO}^{\rm pers}\ge0.1$ and $\Phi_{\rm CR}\ge0.1$). ``Last half'' is the fraction of runs in
  which the median birth time of the persistent isolated class is later than the median birth time of all
  births. A dash means not measured. The O13b entries for \textsf{C1} come from an independent 60-run
  ensemble.}
  \label{tab:calib-dev}
  \small
  \resizebox{\textwidth}{!}{%
  \begin{tabular}{@{}llccc@{}}
    \toprule
    Pattern & & \textsf{C0} ($n=120$) & \textsf{C1} ($n=60$) & \textsf{C2} ($n=60$) \\
    \midrule
    O4 peak below capacity & E & 1.00 [0.97, 1] & 1.00 [0.94, 1] & 1.00 [0.94, 1] \\
    \quad $N_{\rm pk}/K$; empty cells at peak & & 0.42; 0.18 & 0.37; 0.16 & 0.40; 0.13 \\
    O5 natality collapses at high $N$ & E & 1.00 [0.97, 1] & 1.00 [0.94, 1] & 0.98 [0.91, 1] \\
    O6 no rebound & E & 1.00 [0.97, 1] & 1.00 [0.94, 1] & 0.97 [0.89, 0.99] \\
    O7 extinction & E & 1.00 [0.97, 1] & 1.00 [0.94, 1] & 0.97 [0.89, 0.99] \\
    O10 deterioration before peak & E & 1.00 [0.97, 1] & 1.00 [0.94, 1] & 1.00 [0.94, 1] \\
    O11 crowded-born cohorts fail & E & 1.00 [0.97, 1] & 1.00 [0.94, 1] & 1.00 [0.94, 1] \\
    O12 transplant: $P_T(\mathrm{B})$ & E & 0.82 [0.73, 0.88] & 0.83 [0.72, 0.91] & 0.87 [0.76, 0.93] \\
    \quad $P_T(\mathrm{D})$ & & 0/100 [0, 0.04] & 0/60 [0, 0.06] & 0/60 [0, 0.06] \\
    O13 two deteriorated classes & E & 0.98 [0.93, 0.99] & 1.00 [0.94, 1] & 1.00 [0.94, 1] \\
    \quad persistent O13 & & 0 ($\Phi^{\rm pers}_{\rm BO}=0.006$) & 1.00 [0.94, 1] & 1.00 [0.94, 1] \\
    \textbf{O17 Calhoun-like runs} & E & \textbf{0.98 [0.93, 0.99]} & \textbf{1.00 [0.94, 1]} & \textbf{0.95 [0.86, 0.98]} \\
    \midrule
    O8 senescent decline & D & 0.75 [0.67, 0.82] & 0.72 [0.59, 0.82] & 0.00 [0, 0.06] \\
    O8b slow initial decline & D & 1.00 [0.97, 1] & 1.00 [0.94, 1] & 0.92 [0.82, 0.96] \\
    O9 old, female-biased remnant & D & 0.35 [0.27, 0.44] & 0.57 [0.44, 0.68] & 0.17 [0.09, 0.28] \\
    O11b few late-born females conceive & D & 0.99 [0.95, 1] & 0.96 [0.88, 0.99] & 1.00 [0.94, 1] \\
    O12b mixed transplant & D & 0/200 [0, 0.02] & 0/120 [0, 0.03] & 0.07 [0.03, 0.13] \\
    O13b isolated class is later-born & D & -- & 0 [0, 0.06] & 0 [0, 0.06] (MARG 60/60) \\
    \quad $\Delta b/t_{\rm pk}$; last half & & -- & $-0.02$; 0/60 & $+0.13$; 58/60 [0.89, 0.99] \\
    O14 aggregation with free space & D & 0.95 [0.90, 0.98] & 1.00 [0.94, 1] & 1.00 [0.94, 1] \\
    O15 infant mortality $\to$ 100\% & D & 0.02 [0.00, 0.06] & 0 [0, 0.06] & 0.03 [0.01, 0.11] \\
    O15b conceptions without viable young & D & 0.33 [0.25, 0.41] & 0.58 [0.46, 0.70] & 0.63 [0.50, 0.74] \\
    \midrule
    A1--A4 anti-patterns triggered & & 0/120 [0, 0.03] & 0/60 [0, 0.06] & A2: 2/60; others 0 \\
    \bottomrule
  \end{tabular}}
\end{table}

\paragraph{Magnitudes in the development ensemble of \textsf{C0}.}
The candidates also approach several magnitudes of Universe~25, none of which is a test (Section~\ref{sec:circularity}).
Values are medians over the 120 runs of
\textsf{C0}, with bootstrap intervals; the corresponding figure for Universe~25 is in parentheses.
\begin{itemize}
  \item The peak lies at $N_{\rm pk}/K_\gamma=0.416$ [0.411, 0.420] of the damage-threshold capacity
    $K_\gamma=m_c|G|$, which is not a physical capacity (Universe~25: 2200/3840 $=0.57$, comparable only in sign).
  \item At the peak 18\% of the cells are empty (20\% of nest boxes).
  \item The population is still at $\nu_{8b}=0.998$ of its peak $0.31\,t_{\rm pk}$ later (0.93). This is the
    plateau, a consequence of the absence of mortality before senescence.
  \item The decline lasts $(T_{\rm ext}-t_{\rm pk})/t_{\rm pk}=2.13$ [2.00, 2.18] (2.2 projected, $\ge1.84$
    observed). It reaches at least the observed value in 78\% [69, 84] of the runs. The agreement is not robust:
    the ratio depends on where $t_{\rm pk}$ falls on the plateau (Appendix~\ref{app:calib-evs}).
\end{itemize}
The adult social state falls below $0.5$ at $0.28\,t_{\rm pk}$, long before the peak. In the $\alpha=0$
control the clustering of births is weaker: the Gini coefficient of births over $4\times4$ blocks is 0.075,
against 0.168 in \textsf{C0} ($p=2\cdot10^{-17}$; O14).

\subsection{Ablations of \textsf{C0}}
\label{app:calib-ablations}

Table~\ref{tab:calib-ablations} removes the rules of \textsf{C0} one at a time. Without refuges, each of the
three rules is necessary.
\begin{itemize}
  \item \emph{R2}, switched off (full inheritance and individual recovery at $r=0.01$). Without it,
    neither O11 nor O13 is ever reproduced. Viable pups come from the healthiest mothers and inherit their
    state, so the late cohorts are competent.
  \item \emph{AV}. Without it the population collapses into a few huge pools that leave 76\% of the
    lattice empty, which is the [$\ne$ORIG] picture of ``half the colony empty'', and the isolated class
    disappears.
  \item \emph{R1}. Without it, adults recover when the population thins, and extinction (O7) and O13 drop
    to 0.75 and 0.50 at $\alpha=4$. Removing R1 is tolerable only at $\alpha=6$, with lower robustness.
\end{itemize}
The long-lived control (last row; no rules, $\alpha=1.4$) fails O4: the peak reaches 0.83 of the capacity.
They also fail O11 and O13.

\begin{table}[t]
  \centering
  \caption{Ablations of \textsf{C0} (40 replicates per row; $N_0=400$, longevity preset). Fractions of PASS
  runs; O17 with its 95\% Wilson interval. $\rho_{\rm pk}=N_{\rm pk}/K$; $f_0$ = fraction of empty cells at
  the peak; $\Phi_{\rm BO}$, $\Phi_{\rm CR}$ = isolated and crowded deteriorated adults (maximum over
  $t_{\rm pk}$, $t_{\rm pk}+10$, $t_{\rm pk}+20$).}
  \label{tab:calib-ablations}
  \small
  \resizebox{\textwidth}{!}{%
  \begin{tabular}{@{}lcccccccc@{}}
    \toprule
    Variant & O4 & O7 & O11 & O13 & O17 & $\rho_{\rm pk}$ & $f_0$ & $\Phi_{\rm BO}$ / $\Phi_{\rm CR}$ \\
    \midrule
    \textsf{C0} ($\alpha=4$, $\kappa=0.15$, R1, R2) & 1.00 & 1.00 & 0.98 & 0.98 & 0.95 [0.83, 0.99] & 0.42 & 0.19 & 0.12 / 0.15 \\
    \textsf{C0}, $\alpha=6$ & 1.00 & 1.00 & 0.95 & 1.00 & 0.95 [0.83, 0.99] & 0.35 & 0.32 & 0.13 / 0.17 \\
    \textsf{C0}, $\kappa=0.10$ & 1.00 & 1.00 & 0.95 & 0.98 & 0.93 [0.80, 0.97] & 0.37 & 0.33 & 0.11 / 0.31 \\
    \textsf{C0} without R2 & 1.00 & 0.95 & 0.00 & 0.05 & 0.00 [0, 0.09] & 0.51 & 0.12 & 0.08 / 0.18 \\
    \textsf{C0} without AV & 1.00 & 1.00 & 0.68 & 0.00 & 0.00 [0, 0.09] & 0.32 & 0.76 & 0.05 / 0.87 \\
    \textsf{C0} without R1 & 1.00 & 0.75 & 1.00 & 0.50 & 0.40 [0.26, 0.55] & 0.46 & 0.16 & 0.10 / 0.18 \\
    \textsf{C0}, $\alpha=6$, without R1 & 1.00 & 0.98 & 1.00 & 0.85 & 0.82 [0.68, 0.91] & 0.38 & 0.30 & 0.11 / 0.21 \\
    \textsf{C0}, no inheritance ($w=0$, $r=0.15$) & 1.00 & 1.00 & 1.00 & 0.68 & 0.68 [0.52, 0.80] & 0.46 & 0.15 & 0.11 / 0.16 \\
    Long-lived control ($\alpha=1.4$) & 0.25 & 1.00 & 0.10 & 0.00 & 0.00 [0, 0.09] & 0.83 & 0.58 & 0.04 / 0.91 \\
    \bottomrule
  \end{tabular}}
\end{table}

\paragraph{Crowding is the trigger.}
Setting $\gamma=0$ in \textsf{C0} removes the collapse altogether: the population grows past
$2K_\gamma$ and the mean social state stays near 0.42; among adults it stays near 0.6 (\textsf{C1} with $\gamma=0$),
since the mean over all animals includes the pups, born with $s\le0.2$. The failure of socialisation under R2
therefore needs crowding damage to start. Calhoun's emphasis on the social, rather than purely numerical,
consequences of density is consistent with this division of labour~\cite{Calhoun1973}.

\paragraph{The transplant control constrains the attraction.}
With $\alpha=6$ the control transplants (healthy phase-B animals) fail to found a colony in 60\% of the
trials: strong aggregation of eight animals in one cell deteriorates them. This makes O12 not applicable,
and is why $\alpha=4$ was preferred to $\alpha=6$.

\subsection{Local sensitivity of \textsf{C0}}
\label{sec:calib-sensitivity}

Table~\ref{tab:calib-oat} varies each parameter of \textsf{C0} in turn. O4, O5, O6, O7, O10 and O11 remain
at $\ge0.93$ in every variation. The bottleneck is O13: the isolated fraction sits only about 0.02 above
its threshold of 0.1. Two variations break it:
\begin{itemize}
  \item faster learning ($r=0.15$), which shortens the time the young spend deteriorated;
  \item a weaker selectivity of the aversion ($p=2$), which also repels moderately healthy individuals.
\end{itemize}
Both homogenise the occupancy. \textsf{C0} therefore lies inside the robust region but close to its
border in $r$ and $\kappa$. This is consistent with the narrow band of Fig.~\ref{fig:model-schematic}(c),
where $n^*=1/\kappa-1/\alpha$ must stay between about 5.5 and 10.

With refuges the robust region is wider. In the production planes around \textsf{C1}
(Section~\ref{sec:res-robustness}):
\begin{itemize}
  \item the primary outcome stays at $\ge0.83$ for $w\le0.3$ and $r\le0.10$, and falls through O11 at $w\ge0.5$
    and $r\ge0.2$;
  \item it stays at $\ge0.8$ at 34 of the 48 points of the $(\alpha,\kappa)$ plane with $\alpha\ge2$, and fails
    where the free space required by the primary O4 is lost.
\end{itemize}
This local robustness contrasts with the small Calhoun-like fraction of the global Latin hypercube, 2.0\% of the
points with the primary evaluator. That fraction is dominated by the refuge geometry and by thresholds that also
enter the definitions of the classes, and it depends on the sampling box.

\begin{table}[t]
  \centering
  \caption{One-at-a-time sensitivity of \textsf{C0} (30 replicates per row). O17 = fraction of
  Calhoun-like runs with 95\% Wilson interval. In every row the other essential patterns stay at $\ge0.93$.}
  \label{tab:calib-oat}
  \small
  \begin{tabular}{@{}lcc@{\qquad}lcc@{}}
    \toprule
    Variation & O13 & O17 & Variation & O13 & O17 \\
    \midrule
    \textsf{C0} & 0.97 & 0.97 [0.83, 0.99] & $a_{\max}=90$ & 1.00 & 0.93 [0.79, 0.98] \\
    $\alpha=3$ & 0.70 & 0.70 [0.52, 0.83] & $a_{\max}=150$ & 0.97 & 0.93 [0.79, 0.98] \\
    $\alpha=5$ & 1.00 & 1.00 [0.89, 1] & $\gamma=0.07$ & 0.80 & 0.80 [0.63, 0.90] \\
    $\kappa=0.125$ & 1.00 & 0.97 [0.83, 0.99] & $\gamma=0.15$ & 0.97 & 0.90 [0.74, 0.97] \\
    $\kappa=0.175$ & 0.70 & 0.67 [0.49, 0.81] & $w=0.1$ & 0.97 & 0.97 [0.83, 0.99] \\
    $p=2$ & 0.10 & 0.10 [0.03, 0.26] & $w=0.3$ & 0.90 & 0.87 [0.70, 0.95] \\
    $p=8$ & 0.93 & 0.93 [0.79, 0.98] & $a_w=8$ & 0.97 & 0.93 [0.79, 0.98] \\
    $r=0.07$ & 1.00 & 1.00 [0.89, 1] & $a_w=16$ & 0.87 & 0.87 [0.70, 0.95] \\
    $r=0.15$ & 0.17 & 0.17 [0.07, 0.34] & & & \\
    \bottomrule
  \end{tabular}
\end{table}

\subsection{Deteriorated classes: persistence and cohort structure}
\label{sec:calib-classes}

\paragraph{\textsf{C0} reproduces O13 only instantaneously.}
In \textsf{C0} the criterion O13 is met: 12\% of the adults are isolated ($m\le2$) and 14\% crowded
($m>m_c$). But the two groups are snapshots of a single, homogeneously deteriorated population whose
occupancy is overdispersed.
\begin{itemize}
  \item Only 0.6\% of the adults stay isolated for four consecutive weeks.
  \item The weekly autocorrelation of the occupancy experienced by an individual is 0.23.
  \item Isolated and crowded adults do not differ in age (46 versus 46 weeks), sex, social state or total
    affinity (132 versus 145).
\end{itemize}
Changing the affinity memory or formation rate, adding male territoriality, strengthening the aversion
or restricting bond formation to competent pairs does not produce persistence (four-week isolated
fraction $\le0.012$ in every case). At a density of three to four animals per cell every individual
touches neighbours almost every week, and no region is empty enough to stay in.

\paragraph{Refuges.}
Refuge cells of limited capacity, without a preference, remain essentially empty: they hold at most 5\%
of the population and leave the four-week isolated fraction unchanged. A deteriorated animal surrounded by
acquaintances always prefers an occupied neighbouring cell.

With a preference, two conditions are needed:
\begin{itemize}
  \item \emph{The capacity must be hard.} With soft capacity, simultaneous moves produce stampedes into
    the same refuge: the mean crowding reaches 30--60 and the isolated class shrinks.
  \item \emph{The refuges must hold enough animals.} The persistent isolated fraction is bounded by the
    refuge capacity relative to the adult population, $\phi_R=f_Rc_R|G|/N_{\rm ad}$. A class of 10\%
    therefore needs $f_Rc_R\gtrsim0.3$.
\end{itemize}
With $f_R=0.3$, $c_R=2$ and $\pi=100$ (\textsf{C1}), 15\% of the adults form a persistent isolated
class. 95\% of them are in refuges and they remain isolated for 15--16 consecutive weeks on average.
They receive no crowding damage and no social mortality, which matches the ``physically healthy''
beautiful ones of~\cite{Calhoun1973}. Persistent O13 holds in 60/60 runs, and O17 in 60/60
(Table~\ref{tab:calib-dev}).

\paragraph{Cohort structure.}
In \textsf{C1} the persistent isolated animals are of the same age as the crowded ones (51 versus 51
weeks at the peak). They are refugees of the early boom, not a late, unsocialised cohort, and the
pattern O13b fails in every run ($\Delta b=-0.02\,t_{\rm pk}$). Between $t_{\rm pk}$ and $t_{\rm pk}+20$, 17\% of them
are founders, 58--60\% belong to the first generation and 22--27\% to later ones. Most entered the refuges as pups:
62--67\% first entered a refuge before the age of 12 weeks, and 15--20\% before 3 weeks (8 exploratory runs, with
the single-pass and with the exact hard capacity). Pups are independent agents from birth and, being born with $s\le0.2$, are
strongly attracted to refuges; a nest period, in which pups stay in their mother's cell, is a reported
variant (Section~\ref{sec:res-robustness}).

Basing the retreat on the competence reached at the end of the critical window does not help. Under R2
almost every cohort after the first ends the window with $s<0.1$, so there is no gradient between
cohorts. Basing it on the competence of the mother at birth (\textsf{C2}) does create the gradient.
\begin{itemize}
  \item The persistent isolated class is 11 weeks younger than the crowded class (51 versus 63 weeks).
  \item It has 28\% fewer bonds (total affinity 97 versus 136) and a slightly higher social state.
  \item It belongs to the second half of all births in 58/60 runs.
  \item $\Delta b=+0.13\,t_{\rm pk}$, Mann--Whitney $p\approx10^{-46}$.
\end{itemize}
This is the structure Calhoun describes, in which the beautiful ones came from ``the last half of the
population born''~\cite{Calhoun1973}. The price is a broader age distribution in phase~D, which breaks
the senescent rise of the mean age required by O8 (Table~\ref{tab:calib-dev}).

\paragraph{Why O13b cannot pass with its current threshold.}
O13b also requires that half of the isolated class be born after $0.8\,t_{\rm pk}$. With $N_0=400$ this
is out of reach for any retreat mechanism, because of when births happen.
\begin{itemize}
  \item The median birth time of the whole population is $0.18$--$0.24\,t_{\rm pk}$.
  \item Ninety percent of all births occur before $0.44\,t_{\rm pk}$.
  \item Only about 16 pups per run, 0.6\% of all births, are born after $0.8\,t_{\rm pk}$.
\end{itemize}
In Universe~25, by contrast, growth continued almost to the peak, and about half of all births occurred
after $\approx0.8\,t_{\rm pk}$ (our estimate from the growth curve: $N$ reached half its peak around day 435). In the
model the ``peak'' is the top of a plateau: natality dies at
0.5--0.7\,$t_{\rm pk}$, and the colony, made of young adults, hardly changes until senescence. We
therefore also report the cohort criterion relative to the distribution of births (``last half'' in
Table~\ref{tab:calib-dev}). Under that reading \textsf{C2} passes and \textsf{C1} fails.

\subsection{The founder protocol}
\label{sec:calib-founders}

Calhoun founded Universe~25 with four pairs of 48-day-old mice. We tested the corresponding protocol: four
plus four individuals of age 7 weeks placed in the central cell. Table~\ref{tab:calib-founders}
summarises the results.

\paragraph{\textsf{C0} with founders.}
The colony stays within about 10\% of the lattice and deteriorates locally, giving a trivial peak of
$\rho_{\rm pk}\approx0.02$--$0.03$. Under weaker aggregation ($\alpha=1.4$, no R2) the colony grows as a
slow front, at about $0.06$ cells per week. It takes $t_{\rm pk}\approx170$--$215$ weeks to peak, against
a local deterioration time of 20--30 weeks: a Damköhler number $\mathrm{Da}\approx10$.

\paragraph{More mobility.}
Several Moore steps per week (\texttt{move\_substeps}) make things worse. Founders and their offspring
disperse, males and females no longer meet in the same cell, and the colony fails by an encounter-limited
Allee effect. With $N_0=400$, more mixing also homogenises the occupancy and destroys O13: O17 equals
0.93, 0.80, 0.30 and 0 for 1, 2, 4 and 8 sub-steps.

\paragraph{Mating in the Moore neighbourhood.}
With \texttt{mate\_radius}${}=1$, two sub-steps and founders whose initial social state is reduced to
$0.3$ (a stand-in for the adaptation after transfer, phase~A), O1--O3 become applicable and are
frequently reproduced. However, the protocol is not yet robust:
\begin{itemize}
  \item O17 is only 0.17 [0.09, 0.28];
  \item O5 and O13 fail in a third to a half of the runs;
  \item the peak comes too late relative to the life span ($t_{\rm pk}/a_{\max}\approx1.4$, against 0.7
    in Universe~25), so the senescent decline is too short and O8 fails.
\end{itemize}
The production plane of this protocol (initial social state $\times$ sub-steps $\times$ $a_{\max}$, 30 runs
per point, Appendix~\ref{app:res-p4p6}) does not change the conclusion. Its best point ($s_0=0.2$, one
sub-step, $a_{\max}=200$) reproduces O1--O3 jointly with O17p in 0.43 [0.27, 0.61] of the runs, and with the preset
longevity $a_{\max}=120$ no point exceeds 0.13.

\begin{table}[t]
  \centering
  \caption{Founder protocol (four pairs of 7-week-old animals in one cell). Fractions of PASS runs, with
  95\% Wilson intervals where $n\ge30$; counts for the exploratory row with $n=8$. ``No latency'' founders
  start with $s=1$; ``adaptation'' founders start with $s=0.3$. Same-cell rows use $s=1$.}
  \label{tab:calib-founders}
  \small
  \resizebox{\textwidth}{!}{%
  \begin{tabular}{@{}lcccccccc@{}}
    \toprule
    Variant & $n$ & $\rho_{\rm pk}$ & O1 & O2 & O3 & O5 & O13 & O17 \\
    \midrule
    \textsf{C0}, same-cell mating, 1 sub-step & 30 & 0.02 & 0.43 & 0.33 & 1.00 & 0.83 & 0.55 & 0.13 [0.05, 0.30] \\
    \textsf{C0}, same-cell mating, 4 sub-steps & 30 & 0.01 & 0.10 & 1.00 & 1.00 & 0.33 & 0.00 & 0 [0, 0.11] \\
    \textsf{C1}, Moore mating, 2 sub-steps, no latency & 8 & 0.27 & 0/8 & 8/8 & 8/8 & 5/8 & 5/8 & 0/8 \\
    \textsf{C1}, Moore mating, 2 sub-steps, adaptation & 60 & 0.24 & 0.58 [0.46, 0.70] & 0.97 [0.88, 0.99] & 0.71 [0.59, 0.81] & 0.65 [0.52, 0.76] & 0.54 [0.42, 0.66] & 0.17 [0.09, 0.28] \\
    \bottomrule
  \end{tabular}}
\end{table}

\subsection{Generations and chronology}
\label{app:calib-generations}

Table~\ref{tab:calib-generations} follows every animal of \textsf{C1} by generation: the founders are F0, and a pup
belongs to the generation of its mother plus one. The realised fecundity is counted over the whole life of each
female. The table uses the 200 runs of the reference ensemble: the same parameters and seeds were run
again with a register of generations, and every run reproduces the trajectory of the ensemble
(\texttt{scripts/analysis/generaciones\_base.py}). Table~\ref{tab:calib-gen-variants} repeats the summary for the
other members of the reference ensemble and for the variants of Section~\ref{sec:res-robustness}, re-running the first
20 replicates of each production arm with the same seeds (\texttt{scripts/analysis/gamma\_v3.py}; every replicate
reproduces the maximum population of its run). Table~\ref{tab:calib-gamma} gives the deterioration and generation
times, and Table~\ref{tab:calib-chrono} compares the chronology with Universe~25.

\begin{table}[t]
  \centering
  \caption{Generations in \textsf{C1} (reference ensemble, 200 runs). Births: mean over runs and share
  of all births. $s_{\rm nb}$: social state at birth; $s(a_w)$: at the end of the critical window, after which $s$
  can only decrease; $R_{\rm net}$: daughters per female of that generation; ``bred'': fraction of females with at
  least one surviving pup. Founders are 400 animals of age 4 weeks with $s=1$. Later generations add 0.1\% of the
  births.}
  \label{tab:calib-generations}
  \small
  \resizebox{\textwidth}{!}{%
  \begin{tabular}{@{}lcccccccc@{}}
    \toprule
    Generation & Births & Share & Median birth week & $s_{\rm nb}$ & $s(a_w)$ & $s(a_w)\ge0.5$ & Pups per female & $R_{\rm net}$ (bred) \\
    \midrule
    F0 & 400 & -- & -- & 1 & 0.90 & 91\% & 7.6 & 3.82 (73\%) \\
    F1 & 1529 & 69.2\% & 10 & 0.19 & 0.23 & 20\% & 0.80 & 0.40 (16\%) \\
    F2 & 611 & 27.6\% & 19 & 0.14 & 0.13 & 6\% & 0.22 & 0.11 (6\%) \\
    F3 & 67 & 3.0\% & 28 & 0.10 & 0.11 & 2\% & 0.11 & 0.05 (4\%) \\
    F4 & 4 & 0.2\% & 38 & 0.07 & 0.08 & 0 & 0.06 & 0.03 (2\%) \\
    \bottomrule
  \end{tabular}}
\end{table}

\paragraph{A boom of one to two generations.}
The founders deteriorate by crowding during their own adult life: they reach the end of the window with $s=0.90$,
but have $s=0.20$ at 26 weeks of age (week 22). They produce 69\% of all births, about 6\% of their potential
fecundity. Their offspring are born at $s_{\rm nb}=w\,s_f\simeq0.19$, below the deterioration threshold, and gain only
$+0.04$ by the end of the window, because their teachers are already deteriorated; without crowding damage the mean
adult state stays near 0.6. The net reproductive number falls by a factor of about ten in a single step, from 3.8 to
0.40, and the third and later generations add 3.2\% of the births. The generational gradient is monotone
($s_{\rm nb}$ 1, 0.19, 0.14, 0.10; $R_{\rm net}$ 3.8, 0.40, 0.11, 0.05), but it is almost complete between the
founders and the first generation. The structure is the same in the other members of the reference ensemble and in
every variant of Table~\ref{tab:calib-gen-variants}: the mean generation of the births is 1.2--1.6, and the first
generation with $R_{\rm net}<1$ is always F1. It is therefore not an artefact of the synchronous cohort. The
long-lived control without new rules has a \emph{deeper} boom (30 replicates of the null ensemble): 22\% of the births
belong to the third and later generations, the mean generation is 1.95 [1.90, 1.99] against 1.34, and $R_{\rm net}$ of
F1 is 1.21 [1.17, 1.25] against 0.40 (medians with bootstrap intervals; every run of the control has $R_{\rm net}>1$ in
F1, none of \textsf{C1}). The long-lived control also differs from \textsf{C1} in $\alpha$, so the clean contrast is
\textsf{C1} without its four rules (same $\alpha=4$ and crowding damage): mean generation 1.59 and $R_{\rm net}$ of F1
0.66, against 1.34 and 0.40 (Table~\ref{tab:calib-gamma}). The rules, R2 in particular ($-$R2$^*$: 1.62), shorten the
generational depth, because pups no longer inherit the competence of the few healthy mothers.

\paragraph{Deterioration time and generation time.}
Table~\ref{tab:calib-gamma} measures how long deterioration takes. Three means of $s$ are followed every week: over the
adults (age $\ge a_{\min}=6$, the S$_{\rm adult}$ of O10), over the animals past the critical window (age $\ge a_w=12$)
and over the founders. The adult mean is diluted: pups of the first generation are born with $s\le0.2$ and count as
adults from the age of 6 weeks, while they are still inside the window. It falls below 0.5 in week 13, when the
founders still have $\bar s=0.66$ and are 39\% of the adults. The animals past the window, which no longer learn and
admit newborns only 12 weeks after their birth, and the founders, a closed cohort, cross 0.5 in week 16. We therefore
define $t_{\rm det}=t_{1/2}-t_{\rm on}$ on the animals past the window: $t_{\rm det}\approx10$ weeks and
$\Gamma=t_{\rm det}/T_{\rm gen}\approx1.1$ in \textsf{C1} (interquartile range over runs 1.0--1.29; founders 1.2),
against 0.8 with the adult mean of the primary evaluator. Deterioration takes of the order of one generation time. The value
depends on the definition of $T_{\rm gen}$: with overlapping cohorts the interval between median birth weeks is not a
demographic generation time, and in the long-lived and rule-free null models it drops to 5 weeks (4 in some runs),
below the age of first reproduction, which inflates their $\Gamma$. With the standard definition, the mean age of the
mothers at birth ($T_{\rm gen}=14.3$ weeks in \textsf{C1}), $\Gamma\approx0.7$ (interquartile range 0.63--0.74) in
\textsf{C1}, 0.60--0.66 in \textsf{C1}$'$, $-$R2$^*$ and the long-lived and rule-free null models, and 0.85 at
$\gamma=0.02$ (30 and 20 runs, from the same per-run data); $\Gamma$ is of order one under either definition, and
it does not order the depth of the boom under either.

$\Gamma$ is not what keeps the boom shallow. Lowering the crowding damage from $\gamma=0.1$ to 0.02 raises $\Gamma$ from
1.2 to 1.6 (founders 1.2 to 1.8), but the mean generation of the births moves only from 1.35 to 1.42 and $R_{\rm net}$
of F1 from 0.39 to 0.52; the colonies still collapse and pass O17p in 17 of 20 runs. Pups born at
$s_{\rm nb}\approx0.4$ ($s_0=0.25$) change the depth more at about the same $\Gamma$: 1.51 and $R_{\rm net}=0.62$ at
$\gamma=0.1$, 1.64 and 0.84 at $\gamma=0.02$. Across the controls $\Gamma$ does not order the depth either: $-$R2$^*$, with
full inheritance, has $\Gamma=1.2$ and a mean generation of 1.62; $\alpha=0$ has 2.4 and 1.57; the long-lived control
1.6 and 1.95. The depth is set in the first
step, F0$\to$F1: pups are born below the deterioration threshold, conception ($\propto s_fs_m$) and pup survival
($\propto s_f$) fall with the parents' state, and, with R2, pups learn from neighbourhoods made mostly of other pups and
of founders that are already deteriorating. $\Gamma$ remains a heuristic for the comparison with Universe~25
(Section~\ref{sec:model-groups}).

\begin{table}[t]
  \centering
  \caption{Deterioration time and generation time (medians over runs; $n$ runs). $t_{\rm on}$: first week with 5\% of
  the individuals in cells with $m>m_c$. $t_{1/2}$: first week with mean $s<0.5$ over the adults (age $\ge a_{\min}=6$,
  as in O10), over the animals past the critical window (age $\ge a_w=12$, the definition used in the text) and over
  the founders (F0). $T_{\rm gen}$: interval between the median birth weeks of F1 and F2.
  $\Gamma=(t_{1/2}-t_{\rm on})/T_{\rm gen}$, with medians of numerator and denominator; with the mean age of the
  mothers at birth as $T_{\rm gen}$ (14.3 weeks in \textsf{C1}) $\Gamma\approx0.7$, 0.60--0.66 in \textsf{C1}$'$, $-$R2$^*$,
  long-lived and no rules, and 0.85 at $\gamma=0.02$, $s_0=0$, so $\Gamma$ is of order one under either definition. $\langle g\rangle_B$: mean
  generation of the births. Upper block: replicates of the reference ensemble and of the null ensemble (same seeds;
  $-$R2$^*$: $w=1$ without social learning; none: \textsf{C1} without its four rules). Lower block: \textsf{C1} with
  crowding damage $\gamma$ and newborn base $s_0$ varied (other seeds). Script: \texttt{scripts/analysis/gamma\_v3.py}.}
  \label{tab:calib-gamma}
  \small
  \resizebox{\textwidth}{!}{%
  \begin{tabular}{@{}lccccccc@{}}
    \toprule
    Variant & $n$ & $t_{\rm on}$ & $t_{1/2}$: $a\ge a_{\min}$ / $a\ge a_w$ / F0 & $T_{\rm gen}$ & $\Gamma$: $a\ge a_{\min}$ / $a\ge a_w$ / F0 & $\langle g\rangle_B$; $R_{\rm net}$ F1 & O17p \\
    \midrule
    \textsf{C1} & 200 & 6 & 13 / 16 / 16 & 9 & 0.8 / 1.1 / 1.2 & 1.34; 0.40 & 0.995 \\
    \textsf{C1}$'$ & 30 & 5.5 & 12 / 15 / 15 & 9 & 0.8 / 1.0 / 1.1 & 1.34; 0.36 & 0.87 \\
    \textsf{C0} & 30 & 8 & 15 / 17 / 18 & 7 & 0.9 / 1.3 / 1.4 & 1.36; 0.45 & 0 \\
    \textsf{C1}, $\alpha=0$ & 30 & 13 & 22 / 24 / 26 & 5 & 1.8 / 2.4 / 2.6 & 1.57; 0.89 & 0 \\
    $-$R2$^*$ (null) & 30 & 7 & 14 / 15 / 15 & 6.5 & 1.1 / 1.2 / 1.2 & 1.62; 0.69 & 0 \\
    Long-lived (null) & 30 & 11 & 19 / 19 / 20 & 5 & 1.6 / 1.6 / 1.8 & 1.95; 1.21 & 0 \\
    No rules (null) & 30 & 8 & 15 / 15 / 15 & 5 & 1.4 / 1.6 / 1.6 & 1.59; 0.66 & 0 \\
    \midrule
    $\gamma=0.1$, $s_0=0$ & 20 & 6 & 13 / 16 / 16 & 8.5 & 0.9 / 1.2 / 1.2 & 1.35; 0.39 & 0.95 \\
    $\gamma=0.05$, $s_0=0$ & 20 & 5 & 14 / 17 / 17.5 & 8 & 1.1 / 1.4 / 1.5 & 1.37; 0.43 & 1.00 \\
    $\gamma=0.02$, $s_0=0$ & 20 & 6 & 15 / 18 / 20 & 8 & 1.3 / 1.6 / 1.8 & 1.42; 0.52 & 0.85 \\
    $\gamma=0.1$, $s_0=0.25$ & 20 & 6 & 14 / 15 / 15 & 7 & 1.1 / 1.3 / 1.3 & 1.51; 0.62 & 0.05 \\
    $\gamma=0.05$, $s_0=0.25$ & 20 & 6 & 14 / 15.5 / 16 & 7 & 1.1 / 1.4 / 1.4 & 1.54; 0.66 & 0 \\
    $\gamma=0.02$, $s_0=0.25$ & 20 & 6 & 16 / 17 / 17.5 & 7 & 1.3 / 1.6 / 1.6 & 1.64; 0.84 & 0 \\
    \bottomrule
  \end{tabular}}
\end{table}

\begin{table}[t]
  \centering
  \caption{Generational structure (medians over runs; $n$ runs). Share of births of F1, F2 and later generations;
  mean generation of the births $\langle g\rangle_B$; net reproductive number of the founders and of F1;
  $t_{1/2}$: first week with mean $s<0.5$ over the adults ($a\ge a_{\min}$) and over the animals past the window
  ($a\ge a_w$; Table~\ref{tab:calib-gamma}); plateau: weeks with $N\ge0.98N_{\rm pk}$. Upper block: the reference
  ensemble ($t_{1/2}$ for $a\ge a_w$ from 30 replicates for \textsf{C1}$'$, \textsf{C0} and $\alpha=0$).
  Lower block: the first 20 replicates of each production arm of Section~\ref{sec:res-robustness} (30 for the
  long-lived null model), re-run with the same seeds and a register of generations; their O17p over $n=200$ is in
  that section.}
  \label{tab:calib-gen-variants}
  \small
  \resizebox{\textwidth}{!}{%
  \begin{tabular}{@{}lccccccc@{}}
    \toprule
    Variant & $n$ & F1 / F2 / F3+ (\%) & $\langle g\rangle_B$ & $R_{\rm net}$ F0 / F1 & $t_{1/2}$: $a\ge a_{\min}$ / $a\ge a_w$ & Plateau \\
    \midrule
    \textsf{C1} & 200 & 69 / 28 / 3.2 & 1.34 & 3.8 / 0.40 & 13 / 16 & 51 \\
    \textsf{C1}$'$ & 200 & 70 / 26 / 4.0 & 1.34 & 3.4 / 0.37 & 12 / 15 & 50 \\
    \textsf{C0} & 200 & 68 / 29 / 3.1 & 1.35 & 4.4 / 0.44 & 15 / 17 & 54 \\
    \textsf{C1}, $\alpha=0$ & 200 & 50 / 44 / 6.2 & 1.56 & 5.9 / 0.87 & 22 / 24 & 59 \\
    \midrule
    \textsf{C1}, initial ages 4--52 weeks & 20 & 68 / 29 / 3.5 & 1.36 & 3.8 / 0.42 & 12 / 14.5 & 28.5 \\
    \textsf{C1}, initial ages 4--80 weeks & 20 & 62 / 33 / 4.9 & 1.43 & 3.4 / 0.52 & 12.5 / 15 & 22 \\
    \textsf{C1}, nest period of 3 weeks & 20 & 79 / 19 / 1.1 & 1.22 & 3.7 / 0.25 & 12 / 15 & 51.5 \\
    \textsf{C1}, $\mu_{\rm bg}=0.003$ per week & 20 & 68 / 28 / 3.3 & 1.36 & 3.8 / 0.41 & 13 / 16 & 19 \\
    \textsf{C1}, $\mu_{\rm bg}=0.008$ per week & 20 & 65 / 30 / 5.0 & 1.40 & 3.7 / 0.48 & 14 / 17 & 11 \\
    \textsf{C1}, asynchronous update & 20 & 67 / 29 / 3.0 & 1.36 & 4.5 / 0.45 & 14 / 16 & 57 \\
    \textsf{C1}, $s_0=0.25$ & 20 & 58 / 34 / 7.6 & 1.51 & 3.7 / 0.60 & 14 / 15 & 52 \\
    \textsf{C1}, $s_0=0.5$ & 20 & 52 / 38 / 11 & 1.60 & 3.6 / 0.72 & 14 / 15 & 54 \\
    Long-lived, no new rules (null) & 30 & 35 / 42 / 22 & 1.95 & 4.9 / 1.21 & 19 / 19 & 45.5 \\
    \bottomrule
  \end{tabular}}
\end{table}

\begin{table}[t]
  \centering
  \caption{Chronology of \textsf{C1} (reference ensemble, 200 runs; medians, interquartile range in
  brackets) and of Universe~25 (day/7; \cite{Calhoun1973}). $t'_{\rm pk}$: first week with $N\ge0.98N_{\rm pk}$.
  ``Functional B'': growth while the animals past the critical window ($a\ge a_w$), which newborns do not dilute,
  keep a mean $s\ge0.5$ (Table~\ref{tab:calib-gamma}); the founders cross 0.5 in the same week. The adult mean
  ($a\ge a_{\min}$, as in O10) crosses 0.5 three weeks earlier because pups count as adults while still inside the
  window.}
  \label{tab:calib-chrono}
  \small
  \begin{tabular}{@{}lcccc@{}}
    \toprule
    Event & \textsf{C1} (week) & $/t_{\rm pk}$ & Universe~25 (week) & $/t_{\rm pk}$ \\
    \midrule
    First litters (end of A) & 3 [0] & 0.06 & 15 & 0.19 \\
    Mean $s<0.5$, adults $a\ge a_{\min}$ & 13 [1] & 0.24 & -- & -- \\
    Mean $s<0.5$, $a\ge a_w$ (end of functional B) & 16 [1] & 0.28 & 45 (end of B) & 0.56 \\
    Half of all births & 13 [1] & 0.23 & $\approx62$ (our estimate) & $\approx0.8$ \\
    90\% / 99\% of all births & 25 [2] / 40 [5] & 0.45 / 0.75 & -- & -- \\
    Start of the plateau, $t'_{\rm pk}$ & 34 [4] & 0.63 & no plateau & -- \\
    $t_{\rm pk}=\arg\max N$ & 55 [10] & 1 & 80 & 1 \\
    Last surviving birth & 59 [16] & 1.00 & 86 & 1.07 \\
    Last conception & 61 [16] & 1.02 & 131 & 1.64 \\
    End of the plateau (first deaths) & 86 [1] & 1.54 & -- & -- \\
    Extinction & 169 [15] & 3.07 & 254 (projected) & 3.18 \\
    \bottomrule
  \end{tabular}
\end{table}

\paragraph{The phases are compressed.}
In terms of Calhoun's phases (Table~\ref{tab:calib-chrono}), the model has a growth with competent adults
(functional B) from week 3 to week 16: 13 weeks, about two doublings (from 400 to about 1900 animals) and two thirds of
all births. In Universe~25 phase~B lasted 30 weeks and about five doublings, and most births came later, in phase~C.
Measured with the adult mean of the primary evaluator, which pups dilute, functional B would end in week 13. Growth with
deteriorated adults (C) lasts from week 16 to the start of the plateau at week 34; then the population neither
grows nor declines for about 51 weeks, and the senescent decline (D) lasts from week 86 to week 169. The sequence
A--B--C--D appears in the right order, but B and C are compressed, natality is front-loaded, and a plateau without
births or deaths, absent from Universe~25, is inserted before D. O10b, which asks for a functional growth phase of
half the time to the peak, passes in 2\% of the runs with the adult mean of the primary evaluator
(Section~\ref{sec:res-production}). Part of that failure is the same dilution that affected the preregistered O10. Measured
on the animals past the window, the functional growth lasts $0.46\,t'_{\rm pk}$ (median), and O10b passes in 24\% of the
200 runs and is MARGINAL in the rest (founders: 38\%). It still falls short of Calhoun's phase~B, which lasted
$0.56\,t_{\rm pk}$ and 30 weeks. Measured on these animals O10 is unchanged: it passes in every run of \textsf{C1} and, as
with the adult mean, in three of the four null models ($-$R2$^*$, long-lived and no rules, 30 replicates each), so it remains
generic. The same holds for the founders' mean.

\paragraph{Non-synchronous initial condition.}
In the production stage ($n=200$, primary evaluator), initial ages drawn uniformly over 4--52 weeks give
O17p $=0.93$ [0.89, 0.96], and ages over 4--80 weeks, which include post-reproductive founders, give 0.63 [0.56,
0.69] (0.68 and 0.31 in \textsf{C1}$'$). Only O5 changes, and it becomes MARGINAL rather than failing: 14 MARGINAL
and no FAIL with 4--52 weeks, 73 MARGINAL and one FAIL with 4--80 weeks. Its level clause is always met:
$\nu_5=N(t_{B99})/N_{\rm pk}$ has a median of 0.997 and a minimum of 0.97 over the 200 runs of 4--80 weeks, so natality
still ends at the population maximum. What fails is the timing clause. The older founders die earlier (plateau of 28
and 22 weeks), their deaths balance the last births, and the plateau starts earlier ($t'_{\rm pk}$ 30 and 29 weeks,
against 34), while the last 1\% of the births still comes at weeks 38--41. Hence
$\Delta_5=(t_{B99}-t'_{\rm pk})/t'_{\rm pk}$ rises from 0.17 to a median of 0.22 and 0.42, and lies between 0.52 and 0.97
in the MARGINAL runs (series of the production ensembles). The classes, O11, the low peak and the generational
structure (Table~\ref{tab:calib-gen-variants}) do not change: O11, O13 and O13p pass in every run of both arms. The
primary outcome therefore depends on the initial condition only through the timing of O5 relative to $t'_{\rm pk}$, a
generic criterion that carries no evidence for the rules (Section~\ref{sec:res-robustness}).

\paragraph{Pups born above the deterioration threshold.}
With $w=0.2$ and $s_0=0$ every pup is born with $s\le0.2$, which makes part of O11 definitional. With $s_0=0.25$ or
$0.5$ pups are born at $s_{\rm nb}\approx0.4$ or $0.6$. In the production stage ($n=200$) extinction,
the low peak ($\rho_{\rm pk}=0.40$ and 0.44), the two classes and O5 still hold, but O11 passes in 8\% and 0\% of the
runs, and so O17p is 0.08 and 0.00 (0.64 and 0.36 with the preregistered O11; Section~\ref{sec:res-robustness}). O11
reads the state at the age of maturity, 6 weeks, when these pups still have $s\approx0.36$--0.47 (exploratory runs).
At the end of the window crowding has brought them down: in 8 exploratory runs per value the last half of the births
reaches $s(a_w)\approx0.12$, and only 7\% of its females breed. The failure of the crowded-born cohorts as O11 measures
it is therefore largely definitional, tied to $s_0=0$; measured at the end of the window it survives, but it is then
produced by crowding damage, as in the long-lived control, and not by birth below the threshold. The boom remains
shallow: the mean generation of the births is 1.51 and 1.60, and $R_{\rm net}$ of F1 is 0.60 and 0.72 (20 replicates
of each arm; Table~\ref{tab:calib-gen-variants}).

\subsection{Scales and the mapping to Universe~25}
\label{app:calib-scales}

Table~\ref{tab:calib-scales} lists the scales of \textsf{C1} with their counterparts in Universe~25 and states which
were calibrated. The mapping is the one of Section~\ref{sec:model-rules}: a cell is a resting site, $m_c$ the
comfortable occupancy, and the refuges are the high, less preferred nest boxes.

\begin{table}[t]
  \centering
  \caption{Scales of \textsf{C1} and of Universe~25 \cite{Calhoun1973}. Last column: calibrated in this work (C),
  taken from Calhoun (U), a modelling choice of the core model (M), or emergent (E). $D_{\rm eff}$ in cells$^2$ per week. $^\dagger$Our estimate. $^\ddagger$$f_R$ and $\pi$
  are calibrated; $c_R=2$ is the isolation threshold of O13, not the capacity of a nest box. $^\S$Used informally in
  the calibration and part of the primary O4 ($f_0\ge0.10$; Section~\ref{sec:circularity}). The model's
  $t_{\rm det}$ is measured on the animals past the critical window (Table~\ref{tab:calib-gamma}); for Universe~25 it is
  the duration of phase~C, a different operational definition, so the two values of $\Gamma$ compare only in order of
  magnitude.}
  \label{tab:calib-scales}
  \small
  \begin{tabular}{@{}llll@{}}
    \toprule
    Quantity & Model (\textsf{C1}) & Universe~25 & \\
    \midrule
    Time step & 1 week & -- & M \\
    Resting sites & 900 cells & 256 nest boxes & M \\
    Comfortable occupancy & $m_c=8$ & 15 per box & M \\
    $K_\gamma$ (sites $\times$ occupancy) & 7200 (5580 with refuges) & 3840 & -- \\
    Isolation sites & 270 refuges, 540 places & high boxes, 25--50\%$^\dagger$ & C$^\ddagger$ \\
    Initial population & 400, age 4 weeks & 8, age 7 weeks & M \\
    Peak; $N_{\rm pk}/K_\gamma$ & 2600; 0.36 & 2200; 0.57 & E \\
    Occupancy per site at the peak & 2.9 & 8.6 & E \\
    Empty sites at the peak & 16\% & 20\% & E$^\S$ \\
    Movement & $\le1$ cell/week; $D_{\rm eff}\approx0.13$ & enclosure crossed daily & C ($\alpha$) \\
    Contact neighbourhood & $3\times3$ cells & cell (0.36\,m$^2$), whole floor & M \\
    Maturity / menopause / life & 6 / 80 / 120 weeks & 6--7 / 80 / 114 weeks & U \\
    Litter; interval between litters & 6.55; $\ge3.3$ weeks & 6--8; 3--4 weeks$^\dagger$ & M \\
    Gestation; weaning & none; none & 3 weeks; 3 weeks & -- \\
    Critical window & 12 weeks & youth & C \\
    Doubling time during growth & 6--6.5 weeks & 4.7--7.9 (B), 21 (C) weeks & E \\
    Generation time $T_{\rm gen}$ & $\approx9$ weeks & $\approx10$ weeks$^\dagger$ & E \\
    Deterioration time $t_{\rm det}$ ($a\ge a_w$) & 10 weeks & $\approx35$ weeks (phase C)$^\dagger$ & E \\
    $\mathrm{Da}=t_{\rm mix}/t_{\rm det}$ & $\approx1.7\times10^2$ & $\lesssim4\times10^{-3}$ & E \\
    $\Gamma=t_{\rm det}/T_{\rm gen}$ & 1.1 & $\approx3.5^\dagger$ & E \\
    \bottomrule
  \end{tabular}
\end{table}

\subsection{The extinction time as an extreme-value statistic}
\label{app:calib-evs}

Without mortality before senescence the colony dies when its last member dies, $T_{\rm ext}=\max_i(b_i+\ell_i)$,
with $b_i$ the birth time and $\ell_i$ the lifetime of animal $i$. Table~\ref{tab:calib-evs} decomposes it as
$T_{\rm ext}=t_{\rm last}+(T_{\rm ext}-t_{\rm last})$, where $t_{\rm last}$ is the last surviving birth.
\begin{itemize}
  \item The second term, the longest residual lifetime, is about 110 weeks in every model, with an interquartile
    range of 5--6 weeks: of the order of the inverse slope of the logistic hazard, $1/k\approx7$ weeks.
  \item The spread of $T_{\rm ext}$ is therefore that of $t_{\rm last}$ (correlation 0.89--0.98). That spread is set
    by how abruptly natality stops, measured by the e-folding time of the tail of births,
    $\tau_B=(t_{B99}-t_{B90})/\ln10$.
  \item The relative spread is small for two trivial reasons: the median contains a nearly deterministic lifetime of
    110 weeks, and the numerator is set by $\tau_B$. Without attraction ($\alpha=0$) natality stops more abruptly
    ($\tau_B=2.6$ weeks) and $T_{\rm ext}$ is \emph{more} reproducible than in \textsf{C1}; without AV (\textsf{C1}$'$)
    the tail of births is longer ($\tau_B=8.7$) and so is the spread.
\end{itemize}

\begin{table}[t]
  \centering
  \caption{Decomposition of the extinction time (reference ensemble, 200 runs per model; medians,
  interquartile range in brackets; one run of \textsf{C1}$'$ is censored).}
  \label{tab:calib-evs}
  \small
  \begin{tabular}{@{}lcccccc@{}}
    \toprule
    Model & $T_{\rm ext}$ & IQR/median & $t_{\rm last}$ & $T_{\rm ext}-t_{\rm last}$ & corr$(T_{\rm ext},t_{\rm last})$ & $\tau_B$ \\
    \midrule
    \textsf{C1} & 169 [15] & 0.089 & 59 [16] & 111 [5] & 0.91 & 6.9 \\
    \textsf{C1}$'$ & 182 [22] & 0.121 & 71 [22] & 110 [6] & 0.98 & 8.7 \\
    \textsf{C0} & 166 [14] & 0.084 & 56 [13] & 111 [5] & 0.92 & 6.1 \\
    \textsf{C1}, $\alpha=0$ & 155 [7] & 0.045 & 45 [10] & 111 [6] & 0.89 & 2.6 \\
    \bottomrule
  \end{tabular}
\end{table}

\paragraph{Dependence on the size of the system.}
Table~\ref{tab:res-size} runs \textsf{C1} at fixed density on lattices from $L=15$ to $L=60$ ($N_0\propto L^2$, 60
runs each; production stage). The median of $T_{\rm ext}$ grows with the logarithm of the size, by about 6
weeks per unit of $\ln L^2$ (from 162 to 184 weeks), and so does $t_{\rm last}$; the residual lifetime
$T_{\rm ext}-t_{\rm last}$ stays at 110--112 weeks. The relative spread is 0.15 for $L\le20$ and about 0.08 for
$L\ge30$; with 60 runs the intervals at $L\ge30$ overlap, and the slow decrease expected for the maximum of many
events is within them (exploratory runs give 0.05 at $L=84$). The value 0.089 of the reference ensemble is therefore a
value for $L=30$. The peak density (0.35--0.37) and the persistent fraction (0.17--0.18) barely change; the empty
fraction falls from 0.18 to 0.15, an effect of the walls, and O17p rises from 0.82 at $L=15$ to 1.00 at $L\ge45$, so
small lattices fail more often. In exploratory runs up to $L=84$ (30 per size, single-pass refuge admission) no colony recovers
(A2 cleared in all 150 runs, up to about 21\,000 animals): we find no healthy pockets that could reseed a large colony.

